\documentclass[12pt]{article}

\usepackage[a4paper, total={6in, 8in}]{geometry}
\usepackage{textcomp}

\begin{document}
\setlength{\parindent}{0pt}

\def \t {\textrightarrow}
\def \v {\vspace{1em}}
\def \bi {\begin{itemize}}
\def \ei {\end{itemize}}
\def \s[#1] {\section*{#1}}
\def \ss[#1] {\subsection*{#1}}
\def \sss[#1] {\subsubsection*{#1}}

\s[Picture of Dorian Gray]
\ss[Preface]
It is regarded as the manifesto of english aestetism \t here there is a series of aphorisms \\
Aphorism = it is a short clever statement (witty) intended to express a general truth \\
Some of these aphorisms are contradictory, some of them are even obscure (its meaning is not clear) \t other are controversial and provoking \\
Main points made in the preface? \\
Aph. 1 \t "The artist is the creator of beautiful things" \t serves the "art of art's sake" concept \\
Aph. 2 \t "To reveal art and to conceal the artist is art's aim" \\
"Those who find ugly meanings ... charming" \t art should not express anything other than beauty \\
"There is no such thing as an immoral book ... that is all" \t art is not about moral ideas \\
Preface in "Lyrical ballads" \t the language was simple: it was used by pure people, and to create a cultural revolution (to make up for the failure of the political one) \\
Caliban is the co-protagonist of the Tempest in Shakespeare \t the monstruos creature \\
"The artist can express anything" \\
"The type of all the arts is the art of the musician" \t explain? \\
"All art is at one .." \\
ASK THIS PART ABOUT UNDERLING THINS, IMPORTANT \\
Point number 7 ?? \\
He has an elitist vision of art \t only people with the cultural means can understand it

\ss[The studio]
The studio is sorrounded by a beautiful garden \t in line with aestetism \t + expresive objects \\
Exotic form of art \t there isn't any condemnation of forms of art that are not western (japanese art mentioned) \\
One night Basil visits Dorian, he shows him the portrait (which has become the portrait of an ugly man) and killed the painter, who he considers to be the cause \\
Academy is too mainstream and vulgar \t the other one no \\
Aphorism said by Lord Henry \\
One night Dorian and Cicil's brother were sitting in the same opium den \t he wants revenge, recognizes him because he gets called the "Prince charming" \t he tries to kill him, but Dorian convinces him he is not doriang because he is too young \\
Basil is smart, so he cannot be handsome \t the opposite is valid for Dorian \t if you are too beautiful or too smart you will have a hard life \t is better to be stupid or ugly \\
Dorian will cause other people to suffer and will suffer himself \t prophecy
\end{document}
